% GENERATED FILE. Edit canonical lessons, then run npm run generate:books.
% canonical-source-hash: fnv1a64:6ea4e4d03883d820
% canonical-lessons: BN-C219-nokol, BN-C219-jibonto, BN-C219-mrito, BN-C219-phaka, BN-C219-sthaniyo

\chapter{Fake, Alive, Dead, Deserted, Local}
\label{ch:bn-fake-alive-dead-deserted-local}
\hlchaptermodality{219}

\begin{chapteropening}
\textbf{By the end of this chapter:} \emph{I can say fake, alive, dead, deserted and local.}

Complete the last lesson of chapter 219: I can say fake, alive, dead, deserted and local.
\end{chapteropening}

\section[\texorpdfstring{\bn{নকল} (nôkôl)}{নকল (nôkôl)}]{\bn{নকল} (nôkôl) — fake}

\label{lesson:BN-C219-nokol}



\begin{quote}
Before the new one: say the Bengali for then, then the Bengali for that is.
\end{quote}

\subsection*{You'll want to know: \bn{নকল}}
\textbf{\bn{নকল}} — \emph{nôkôl} — \textquotedblleft{}fake\textquotedblright{}.

\begin{grammarlens}[title={What you've built}]
One more word for this chapter.
\end{grammarlens}

\subsection*{Your turn}
\begin{itemize}
\raggedright
  \item \emph{Say it:} \emph{nôkôl}
  \item \emph{Say it:} \emph{nôkôl}, once more
  \item \emph{Say it:} say \emph{nôkôl}
  \item \emph{Recall:} say the Bengali for then, then the Bengali for that is, then say \emph{nôkôl} again
\end{itemize}

\begin{tcolorbox}[breakable,colback=teal!4,colframe=teal!35!black,title={Before you move on}]
What does \bn{নকল} mean? (\textquotedblleft{}Fake\textquotedblright{}.) Say it once more.
\end{tcolorbox}
\section[\texorpdfstring{\bn{জীবন্ত} (jībôntô)}{জীবন্ত (jībôntô)}]{\bn{জীবন্ত} (jībôntô) — alive, living}

\label{lesson:BN-C219-jibonto}



\begin{quote}
Before the new one: say the Bengali for that is, then the Bengali for fake.
\end{quote}

\subsection*{You'll want to know: \bn{জীবন্ত}}
\textbf{\bn{জীবন্ত}} — \emph{jībôntô} — \textquotedblleft{}alive, living\textquotedblright{}.

\begin{grammarlens}[title={What you've built}]
One more word for this chapter.
\end{grammarlens}

\subsection*{Your turn}
\begin{itemize}
\raggedright
  \item \emph{Say it:} \emph{jībôntô}
  \item \emph{Say it:} \emph{jībôntô}, once more
  \item \emph{Say it:} say \emph{jībôntô}
  \item \emph{Recall:} say the Bengali for that is, then the Bengali for fake, then say \emph{jībôntô} again
\end{itemize}

\begin{tcolorbox}[breakable,colback=teal!4,colframe=teal!35!black,title={Before you move on}]
What does \bn{জীবন্ত} mean? (\textquotedblleft{}Alive, living\textquotedblright{}.) Say it once more.
\end{tcolorbox}
\section[\texorpdfstring{\bn{মৃত} (mritô)}{মৃত (mritô)}]{\bn{মৃত} (mritô) — dead}

\label{lesson:BN-C219-mrito}



\begin{quote}
Before the new one: say the Bengali for fake, then the Bengali for alive.
\end{quote}

\subsection*{You'll want to know: \bn{মৃত}}
\textbf{\bn{মৃত}} — \emph{mritô} — \textquotedblleft{}dead\textquotedblright{}.

\begin{grammarlens}[title={What you've built}]
One more word for this chapter.
\end{grammarlens}

\subsection*{Your turn}
\begin{itemize}
\raggedright
  \item \emph{Say it:} \emph{mritô}
  \item \emph{Say it:} \emph{mritô}, once more
  \item \emph{Say it:} say \emph{mritô}
  \item \emph{Recall:} say the Bengali for fake, then the Bengali for alive, then say \emph{mritô} again
\end{itemize}

\begin{tcolorbox}[breakable,colback=teal!4,colframe=teal!35!black,title={Before you move on}]
What does \bn{মৃত} mean? (\textquotedblleft{}Dead\textquotedblright{}.) Say it once more.
\end{tcolorbox}
\section[\texorpdfstring{\bn{ফাঁকা} (phā̃kā)}{ফাঁকা (phā̃kā)}]{\bn{ফাঁকা} (phā̃kā) — deserted, empty}

\label{lesson:BN-C219-phaka}



\begin{quote}
Before the new one: say the Bengali for alive, then the Bengali for dead.
\end{quote}

\subsection*{You'll want to know: \bn{ফাঁকা}}
\textbf{\bn{ফাঁকা}} — \emph{phā̃kā} — \textquotedblleft{}deserted, empty\textquotedblright{}.

\begin{grammarlens}[title={What you've built}]
One more word for this chapter.
\end{grammarlens}

\subsection*{Your turn}
\begin{itemize}
\raggedright
  \item \emph{Say it:} \emph{phā̃kā}
  \item \emph{Say it:} \emph{phā̃kā}, once more
  \item \emph{Say it:} say \emph{phā̃kā}
  \item \emph{Recall:} say the Bengali for alive, then the Bengali for dead, then say \emph{phā̃kā} again
\end{itemize}

\begin{tcolorbox}[breakable,colback=teal!4,colframe=teal!35!black,title={Before you move on}]
What does \bn{ফাঁকা} mean? (\textquotedblleft{}Deserted, empty\textquotedblright{}.) Say it once more.
\end{tcolorbox}
\section[\texorpdfstring{\bn{স্থানীয়} (sthānīyô)}{স্থানীয় (sthānīyô)}]{\bn{স্থানীয়} (sthānīyô) — local}

\label{lesson:BN-C219-sthaniyo}



\begin{quote}
Before the new one: say the Bengali for dead, then the Bengali for deserted.
\end{quote}

\subsection*{You'll want to know: \bn{স্থানীয়}}
\textbf{\bn{স্থানীয়}} — \emph{sthānīyô} — \textquotedblleft{}local\textquotedblright{}.

\begin{grammarlens}[title={What you've built}]
One more word for this chapter.
\end{grammarlens}

\subsection*{Your turn}
\begin{itemize}
\raggedright
  \item \emph{Say it:} \emph{sthānīyô}
  \item \emph{Say it:} \emph{sthānīyô}, once more
  \item \emph{Say it:} say \emph{sthānīyô}
  \item \emph{Recall:} say the Bengali for dead, then the Bengali for deserted, then say \emph{sthānīyô} again
\end{itemize}

\begin{tcolorbox}[breakable,colback=teal!4,colframe=teal!35!black,title={Before you move on}]
What does \bn{স্থানীয়} mean? (\textquotedblleft{}Local\textquotedblright{}.) Say it once more.
\end{tcolorbox}
